% TOP_PROBLEM: {"id": 2307063, "title": "Research Problems in Function Theory — Problem 7.63", "category_id": 8, "category": {"id": 8, "name": "analysis", "display_name": "Analysis", "description": "Limits, continuity, calculus, and function theory.", "slug": "analysis", "order_index": 8, "created_at": "2026-07-31T15:26:25.671Z"}, "status": "open"}
% FIRST_POSTED: null
% ENTRY_KIND: statement_only
% Imported from: batch06.tex; UnsolvedMath ID 2307063
% Compile twice: lualatex 2307063.tex
\documentclass[11pt]{article}
\usepackage{fontspec}
\setmainfont{Latin Modern Roman}
\usepackage{amsmath,amssymb,mathtools,unicode-math}
\setmathfont{Latin Modern Math}
\usepackage{xurl,hyperref}
\hypersetup{hidelinks,pdftitle={UnsolvedMath Problem 2307063}}
\newcommand{\sref}[2]{\hyperlink{source:#1:#2}{[S#2]}}
\newcommand{\cataloguescope}[1]{\paragraph{Mathematical question.}#1}
\begin{document}
% UnsolvedMath ID 2307063; source catalogue code AMR-022-7063
\setcounter{section}{2307062}
\section{Research Problems in Function Theory — Problem 7.63}\label{2307063}
{\small\sffamily UnsolvedMath ID 2307063 · Analysis}
\subsection{Definitions and mathematical statement}

\paragraph{Shared notation.}$\mathbb N=\{1,2,\ldots\}$, and $\mathbb Z,\mathbb Q,\mathbb R,\mathbb C$ denote the integers, rationals, reals and complex numbers. Any other conventions are specified locally.


For $n\ge2$, $\alpha>0$ and $R>0$, the Bochner--Riesz operator on Schwartz functions is defined by
\[\widehat{T_R^\alpha f}(\xi)=\big(1-|\xi|^2/R^2\big)_+^\alpha\hat f(\xi),
\qquad\hat f(\xi)=\int_{\mathbb R^n}f(x)e^{-2\pi i x\cdot\xi}\,dx,
\]
where $s_+=\max(s,0)$ and Schwartz functions are smooth functions all of whose derivatives decay faster than every inverse power. For $f\in L^p$, $p=2n/(n+1)$, interpret $T_R^\alpha f=K_R^\alpha*f$, with $K_R^\alpha$ the inverse Fourier transform of the displayed multiplier. The kernel belongs to the conjugate Lebesgue class, so this defines the convolution for such $f$.
\paragraph{Statement recorded in the supplied source.}
For every $f\in L^{2n/(n+1)}(\mathbb R^n)$ and every $\alpha>0$, does $T_R^\alpha f(x)\to f(x)$ for almost every $x$ as $R\to\infty$? Retain also the source's associated norm-convergence question in dimensions $n\ge3$: does $\|T_R^\alpha f-f\|_{2n/(n+1)}\to0$ for every such $f$ and $\alpha$? The original update reports an affirmative almost-everywhere result in all dimensions by Carbery, Rubio de Francia and Vega. \sref{2307063}{2}
\subsection{Short English statement}
Historical question: do these Fourier summation means converge almost everywhere at the specified Lebesgue exponent, and what is the associated norm-convergence behaviour in higher dimensions?
\subsection{Sources}
\begin{description}
\item[\hypertarget{source:2307063:1}{S1}] ULAM.AI and the UnsolvedMath Contributors, UnsolvedMath: A Curated Collection of Open Mathematics Problems, record 2307063 (AMR-022-7063). CC BY 4.0. \url{https://huggingface.co/datasets/ulamai/UnsolvedMath}.
\item[\hypertarget{source:2307063:2}{S2}] Walter K. Hayman and Edward F. Lingham, Research Problems in Function Theory (2018), Problem 7.63, its complete update and relevant preceding definitions. \url{https://arxiv.org/abs/1809.07200}.
\end{description}
\label{end:2307063}
\end{document}
